\documentclass[11pt,a4paper]{article}
\usepackage[margin=21mm,headheight=15pt]{geometry}
\usepackage[T1]{fontenc}
\usepackage{lmodern,amsmath,amssymb,booktabs,tabularx,enumitem,xcolor,fancyhdr}
\usepackage[hidelinks]{hyperref}
\definecolor{accent}{RGB}{0,114,178}
\setlength{\parindent}{0pt}
\setlength{\parskip}{6pt}
\setlist{nosep,leftmargin=*}
\pagestyle{fancy}
\fancyhf{}
\fancyhead[L]{\small Econometrics 25117 | Instructor guide}
\fancyhead[R]{\small 2026--2027}
\fancyfoot[L]{\small Private teaching notes}
\fancyfoot[R]{\thepage}
\newcommand{\heading}[1]{\section*{\color{accent}#1}}
\newcommand{\cue}[2]{\par\noindent\textbf{#1}\quad #2\par}
\setlength{\emergencystretch}{1em}
\newcommand{\E}{\mathbb{E}}
\newcommand{\Var}{\operatorname{Var}}
\newcommand{\Cov}{\operatorname{Cov}}
\begin{document}
\begin{center}{\Large\bfseries Functional forms and interactions}\end{center}
\textbf{Slide route:} Lecture7v2.tex: polynomials, logarithms, and interactions through wrapping up. Chapter 8.
\heading{Session 9 | A slope need not be constant}
\textbf{Outcomes:} Interpret quadratic terms, distinguish log specifications, and calculate a conditional effect in an interaction model.
\heading{100-minute route}
\begin{tabularx}{\textwidth}{@{}rX@{}}\toprule
0--10 & Retrieve coefficient units and joint tests.\\
10--30 & Nonlinear functions and polynomials.\\
30--50 & Linear-log, log-linear, log-log examples.\\
50--55 & Pause.\\
55--80 & Dummy/continuous and continuous/continuous interactions.\\
80--90 & Integrate interpretation with tests.\\
90--100 & Exit exercise and Midterm 2 preparation.\\\bottomrule
\end{tabularx}
\heading{Worked example | Quadratic experience}
\cue{Use with slides}{Lecture7v2 slides 8--13: Nonlinear Functions and Polynomials in $X$. Use the derivative and turning-point calculation on slides 9--13; pause before moving to logs.}
For the synthetic fitted relationship
\[
\widehat{wage}=10+2x-.05x^2,
\qquad \frac{\partial\widehat{wage}}{\partial x}=2-.1x,
\]
the marginal effect is 1 at $x=10$ and 0 at $x=20$. The turning point is 20. The exact fitted change from 10 to 11 is
\[
2-.05(11^2-10^2)=.95,
\]
not exactly the derivative at 10. Only discuss the turning point substantively if it lies within a relevant supported range.
\cue{Board comment}{The model is nonlinear in $x$ but linear in the coefficients. OLS can estimate it after including both $x$ and $x^2$.}
\heading{Log interpretation crib}
\begin{tabularx}{\textwidth}{@{}lX@{}}\toprule
Model & Interpretation of $\beta$\\\midrule
$Y=\alpha+\beta\log X+u$ & 1\% higher $X$: approximately $.01\beta$ units higher $Y$.\\
$\log Y=\alpha+\beta X+u$ & One unit higher $X$: approximately $100\beta$\% higher $Y$.\\
$\log Y=\alpha+\beta\log X+u$ & $\beta$ is the elasticity of the modeled log relationship.\\\bottomrule
\end{tabularx}

\newpage
\heading{Worked examples | Logs and heterogeneous effects}
\cue{Use with slides}{Lecture7v2 slides 14--21 for the log examples and slides 22--39 for interactions. Use the log-outcome calculation on slides 18--21, then use the program-by-experience example on slides 26--38.}
If a log-outcome model has coefficient $.08$, a one-unit increase changes the fitted log index by $.08$. The corresponding multiplicative change is
\[
100(e^{.08}-1)\simeq8.33\%,
\]
versus the 8\% approximation. Exponentiating a fitted conditional log mean does not automatically give the conditional arithmetic mean of $Y$; retransformation requires attention to the error distribution.

For a level-log model with coefficient 20, a 10\% increase in $X$ implies $20\log(1.10)\simeq1.906$ outcome units, compared with the approximation 2.
\heading{Interaction board example}
Let $D$ be a program indicator and $x$ be experience:
\[
\widehat Y=10+2x+3D+1.5(xD).
\]
For $D=0$, the line is $10+2x$; for $D=1$, it is $13+3.5x$. The program-group contrast at experience $x$ is $3+1.5x$, so it equals 9 at $x=4$.
\cue{Ask}{Is 3 the program effect for everyone? No, it is the contrast at $x=0$. Is 1.5 the entire program effect? No, it is the difference in slopes. A causal reading still needs identification assumptions.}
\cue{Test at $x=4$}{Test $\beta_D+4\beta_{xD}=0$, using the covariance of the two estimates. Testing only the main dummy coefficient does not answer this question.}
\cue{Continuous interaction}{For $Y=\beta_0+\beta_1X+\beta_2Z+\beta_3XZ+u$, the $X$ derivative is $\beta_1+\beta_3Z$. State the value of $Z$ whenever reporting it.}
\cue{Exit answers}{A log-log slope of $-.7$ implies roughly a $.7$\% decrease for a 1\% increase in $X$. A quadratic model can be estimated by OLS. An interaction coefficient is not a universal total effect.}
\cue{If short}{Use one example of each functional form and one interaction in depth; assign repetitive interaction plots for self-study. Proposed Midterm 2 coverage is Lectures 5--7, using prerequisites from Lectures 1--4.}
\textbf{Source:} Lecture7v2.tex; synthetic numerical illustrations.

\end{document}
